\documentclass[10pt,a4paper]{amsart}
\usepackage[margin=19mm]{geometry}
\usepackage{amsmath,amssymb,mathtools}
\usepackage[scr=rsfs,cal=euler]{mathalfa}
\usepackage{xfrac}
\usepackage[unicode,hidelinks,bookmarks=false]{hyperref}
\newcommand{\ie}{i.e.\ }
\newcommand{\cf}{cf.\ }
\newcommand{\resp}{respectively\ }
\newcommand{\Subsection}{\subsection}
\providecommand{\nobreakdash}{\nobreak\hbox{-}}
\setlength{\emergencystretch}{2em}
\multlinegap0pt
\usepackage{tikz}
\usepackage{amssymb}
\usepackage[shortlabels]{enumitem}
\usepackage{xcolor}
\newtheorem{proposition}{Proposition}
\usepackage{cleveref}
\crefname{proposition}{Proposition}{Propositions}
\newcommand{\Explicit}{Explicit}
\newcommand\Span{\mathrm{Span}\,}
\newcommand{\bs}{\backslash}
\newcommand{\cplx}{{(K,\partial)} }
\newcommand{\dr}{\partial}
\newcommand{\explicit}{explicit}
\DeclareMathOperator{\g}{\mathtt{g}}
\newcommand{\ols}[1]{\mskip.5\thinmuskip\overline{\mskip-.5\thinmuskip {#1} \mskip-.5\thinmuskip}\mskip.5\thinmuskip} % overline short
\newcommand{\tx}[1]{\text{#1}}
\title{Characterization of the computed homology and cohomology bases}
\author{Yann-Situ Gazull, Aldo Gonzalez-Lorenzo, Alexandra Bac}
\date{Source: arXiv:2509.09350v1 (CC BY 4.0)}
\begin{document}
\maketitle

\noindent\emph{Source and licence.} Characterization of the computed homology and cohomology bases. Version arXiv:2509.09350v1 (CC BY 4.0), distributed by arXiv under the Creative Commons Attribution 4.0 International licence, http://creativecommons.org/licenses/by/4.0/, as recorded in the arXiv licence metadata for this exact version and shown on the abstract page https://arxiv.org/abs/2509.09350v1. Full licence notice: \texttt{licenses/arxiv-2509.09350-CC-BY-4.0.txt}.\par\smallskip

\noindent\textit{Editorial scope.} Counted unit: the article's proposition prop\_elementary-generators (source0.tex lines 674--677). The article's complete original proof of it is delivered with it (source0.tex lines 678--684, one \texttt{proof} environment of 7 lines). No mathematics is written, repaired or re-derived: the statement and the proof are the article's own lines, in the article's own notation, and no retained line is substituted or re-flowed.  Retained uncounted as the source-local context the proof needs: lines 652--666.  Nothing was omitted from any retained span, so the package carries no omission record.  No retained line contains a citation, so the wrapper carries no bibliography and no bibliography entry.  No retained line contains a figure environment, an image-inclusion command or drawing code, so no image is delivered and the package declares no asset dependency.  Editorial formatting only, in addition: an AMS \texttt{amsart} preamble that restates the article's own declarations and macros used by the retained lines, namely the environments \texttt{proposition} on separate counters, together with the standard packages those lines need.  The retained environments print in this document as Proposition 1 in order of appearance, whereas the article prints the same items with the numbering of its own full document.  The complete native source of the article as distributed in the arXiv e-print for this exact version is bundled unchanged as \texttt{sources/184897/source0.tex}.
    \begin{proposition}[\Explicit\ bases equivalent definitions]\label{prop_explicit-bases-equivalence}
        Given a $q$-homology basis $(\g_i)_{i\leq\beta}$ for a complex $\cplx$, the three following properties are equivalent:
        \begin{align*}
        1.\quad & \forall k,\ \ols{\g_k}\bs\bigcup\nolimits_{i \neq k} \ols{\g_i} \neq \emptyset \quad \tx{and}\quad
         \forall J,\ \iota_q^J: H_q\left(K^J\right) \to H_q(K) \ \tx{injective of rank $|J|$}\\
        2.\quad & \forall J\subseteq \{1, \dots, \beta\}, \quad
        \ker\dr \cap K^J_q \subseteq \Span\big((\g_i)_{i\in J}\big)\\
        3.\quad & \forall k,\ \ols{\g_k}\bs\bigcup\nolimits_{i \neq k} \ols{\g_i} \neq \emptyset \quad \tx{and}\quad
        dim \left( H_q\left(K^{\beta}\right) \right) = \beta
        \end{align*}
        %$$
        %\forall k,\ \ols{\g_k}\bs\bigcup_{i \neq k} \ols{\g_i} \neq \emptyset \qquad %\tx{and}\qquad
        %\beta_q\left(K^k\right) = k
        %$$
    \end{proposition}

    \begin{proposition}\label{prop_elementary-generators}
        The generators of an \explicit\ homology basis are \emph{elementary}.
        A cycle is said to be \emph{elementary} if it is not composed of smaller (non-zero) cycles.
    \end{proposition}

    \begin{proof}
% !TEX root = ../../main.tex
Let $(\g_i)_{i\leq \beta}$ be an \explicit\ $q$-homology basis of $(K,\dr)$ and $a$, $b$ be cycles such that $\ols{a}\sqcup \ols{b} = \ols{\g_k}$ (and so $a+b = \g_k$).
Denote $J := \{k\}$.
We get that $a$ and $b$ are in $\ker\dr \cap K^J_q$ so by~\Cref{prop_explicit-bases-equivalence}(2.) they are in $\Span(\g_k) = \{0, \g_k\}$.
Hence, either $a$ or $b$ is $0$, otherwise we would have $\g_k = \g_k + \g_k = 0$.
    \end{proof}

\end{document}
